% Nouveau EGA V (cote 25): archive cote 25, site page 139 (PDF page 140).
% English translation of the checked French annotated master.
% AI-translated by GPT-6 Astra (OpenAI, ChatGPT Pro mode).
\scanpage{25/139}
% The very faint source is retained with all of its checked uncertainty marks.
\hand{\unsure{New} \uline{\unsure{EGA V}}}

\unsure{IV 16.}

\unsure{A new no.:} \uline{\unsure{Singular and supersingular set of a function,}}

\uline{\unsure{and differential criteria}}.

\unsure{This no. will be used in section 20 on hyperplane sections, but its}

\unsure{natural place seems to me to be in section 16.}

\uline{\unsure{Definition}} \unsure{1.} \unsure{Let X be a regular prescheme, f a section of} $\unsure{\mathcal{O}_X}$ \unsure{. A}

\unsure{point} $\unsure{x\in X}$ \unsure{is called a} \uline{\unsure{singular}} \illegible{3} \uline{\unsure{zero}} \unsure{for f}

\unsure{if we have} $\unsure{f_x \in \mathfrak{m}_x^2}$ \unsure{; an} \uline{\unsure{ordinary singular zero}} \unsure{if it is a singular zero}

\unsure{and if in addition the element of} $\unsure{\mathfrak{m}_x^2/\mathfrak{m}_x^3 \subset \mathrm{Sym}(\mathfrak{m}_x/\mathfrak{m}_x^2)}$ \unsure{(well defined),}

\unsure{interpreted as a quadratic form on the dual} $\unsure{t_x}$ \unsure{of} $\unsure{\mathfrak{m}_x/\mathfrak{m}_x^2}$ \unsure{over} $\unsure{k(x)}$,

\unsure{is a nondegenerate form.} \unsure{(A singular zero which is not} \illegible{2}

\unsure{lier is sometimes called an} \uline{\unsure{ordinary singular zero}}\unsure{).}

\uline{\unsure{Remarks}} \unsure{2.} \unsure{If} $\unsure{x\in V(f)}$ \unsure{, then x is a} \var{\inserted{\unsure{non-}}}{\unsure{non-}}{insertion}{ega25_25_139-v0001}\unsure{singular zero of f iff}

$\unsure{f_x \neq 0}$ \unsure{and x is a} \illegible{1} \unsure{nonsingular (= regular) point of} $\unsure{V(f)}$

\unsure{i.e. x is a regular point of} $\unsure{V(f)}$ \unsure{and} $\unsure{V(f)\neq X}$ \unsure{in a neighborhood of}

% The following five source lines are enclosed and crossed out.
\var{\deleted{\unsure{x. One can prove that if x is singular for f, it is}

\unsure{ordinary iff} $\unsure{f_x\neq 0}$ \unsure{and the multiplicity of} $\unsure{\mathcal{O}_{V(f),x}}$ \unsure{is equal to 2}

\unsure{in this case one also sometimes says that x is an} \uline{\unsure{ordinary singular point}}

\unsure{of the prescheme} $\unsure{V(f)}$ \unsure{(i.e.} \unsure{iff} $\unsure{x\in V(f)}$ \unsure{is non} \illegible{3}

$\unsure{f_x\neq 0}$ \unsure{and the multiplicity of} $\unsure{\mathcal{O}_{V(f),x}}$ \unsure{is} \unsure{2)}}}{}{deletion}{ega25_25_139-v0002}

\uline{\unsure{Definition}} \unsure{3.} \unsure{Let X be a prescheme smooth over S, f a}

\unsure{section of} $\unsure{\mathcal{O}_X}$ \unsure{,} $\unsure{x\in V(f)}$ \unsure{. We say that x is a singular zero over}

\uline{\unsure{singular}} \unsure{over S} \unsure{(resp.} \uline{\unsure{generically ordinary singular}}\unsure{) of f over S if}

\unsure{for every extension} \illegible{1} \unsure{of} $\unsure{k(s)}$ \unsure{and every point} $\unsure{x'}$ \unsure{of} $\unsure{X'}$ \unsure{over x}

\unsure{the localization of} $\unsure{x'}$ \unsure{of the corresponding point} \unsure{and of} $\unsure{X_s}$ \unsure{is a}

\unsure{singular (resp. ordinary singular) zero of} $\unsure{f_{s}}$.

\mspage{\hand{139}}
\tnote{25-139-1}{All uncertainties and gaps of this faint draft are retained, including its incomplete and repeated wording. The unresolved French fragment lier is kept literally; its preceding missing letters are not supplied. The later, clearer version has not been used to replace this page's ordinary singular, generically, the duplicated singular over, or the wording of the final localization clause. The inserted non is attached to singular zero in English.}
